\paragraph{A credit sale must create the claim that can later be collected.}
The fixed population above can follow claims already present at opening. A new
sale poses a different problem: the invoice identifier, customer and payment
terms must enter the state when the receivable arises. Increasing only the
aggregate receivable account leaves later collection without an identified
claim. Reserving zero-valued invoices before they exist would instead confuse
absence with an existing claim whose balance happens to be zero.

The additive \texttt{CompanyMemoryTypes/1} declaration distinguishes fixed
records, exact scalar values and finite maps of typed records. An empty map
contains no claims. Once a sale inserts an invoice, each of its actual fields
has an explicit type, including the company and currency context of its
monetary amount. \texttt{CompanyBookTransition/2} checks that these declarations
agree with the current book and its memory schema. It permits insertion of a
complete map entry and replacement of a scalar field. It refuses duplicate
identities, incomplete records, incompatible types and excess capacity. Map
keys remain single identities even when they contain slashes or tildes.
The frozen schema validator does not express the declared capacity, so the
successor checks the typed population again after staging and before returning
the proposed book. Monetary effects, memory changes and postconditions must
all succeed; a failure leaves the supplied book unchanged.

Consider a separate synthetic company opening with cash of USD~100, inventory
of USD~20 and equity of USD~120. It sells two units costing USD~5 each on
2~January 2027 for consideration of USD~30, due on 31~January. The declared sale
creates the named customer's invoice for USD~30 and records its date and sale
event. At the same step it debits receivables and credits revenue by USD~30,
then debits cost of sales and credits inventory by USD~10. There is no cash
movement. A receipt of USD~10 on 15~January debits cash, credits receivables
and reduces only this invoice to USD~20. Receipts identify both the invoice
and its declared customer; neither another customer's label nor a payment
larger than the remainder is accepted.

\texttt{DeclaredReceivableLifecycle/1} requires the sum of invoice remainders
to equal the receivable account before and after each successful command.
It begins with an explicitly empty claim population; the preceding profile
retains its separate opening-claim example. Assessment uses the same declared
strictly-after-due-date convention and contiguous half-open windows. An
assessment on 31~January records zero overdue; one on 28~February records
USD~20. These operations update invoice memory without a journal or cash
flow. Each invoice records whether it has been assessed and at which date.
An invoice created after an earlier assessment therefore remains visibly
unassessed until a subsequent assessment includes it. A receipt preserves
the old snapshot rather than silently recomputing it. Assessment seals its
date against later sales or receipts intended to change that observation.

With no further transactions in 2027, cash is USD~110, receivables USD~20 and
inventory USD~10. Assets and equity are USD~140, profit is USD~20 and operating
cash flow is USD~10. Provisional year-end carry transfers the profit to
retained earnings and preserves every existing invoice field and the
assessment cursor. Collecting the remaining USD~20 in January 2028 raises
cash to USD~130, clears the receivable and produces operating cash flow of
USD~20 in that year. It leaves the prior overdue snapshot historical until
the next declared assessment records zero.

The versioned \texttt{DeclaredReceivableText/1} adapter uses the existing
event, report and carry syntax with explicitly hashed lifecycle imports.
Its checked executions reproduce the structured requests on two different
charts of accounts. An unresolved or refused command retains earlier
successful transactions and reports. These examples establish a bounded
sale-to-collection composition, including identity creation and carry; they
do not supply source-backed opening invoices, revised payment terms,
impairment, tax, arbitrary method authoring or production durability. Those
remain substantive requirements of the wider company simulator.

